\chapter{Biliary Drainage}

\section*{What Is This Test or Procedure?}
Biliary drainage relieves obstruction of the bile ducts by placing a catheter or stent through an endoscopic, percutaneous, or surgical approach.

\section*{Why This Test or Procedure Is Ordered}
It may be required for obstructive jaundice, cholangitis, gallstones, postoperative injury, strictures, or tumors that block bile flow.

\section*{How This Test or Procedure Is Performed}
During ERCP, a stent may be passed through the papilla into the bile duct. If endoscopic access is not possible, interventional radiology can place a transhepatic drain through the skin and liver. The approach depends on anatomy and cause.

\section*{Risks and Follow-up}
Pancreatitis, bleeding, infection, perforation, catheter blockage, and electrolyte or fluid losses are possible. Drain and stent care, laboratory monitoring, and planned exchange or removal are required.

\section*{References}
\begin{enumerate}
  \item MedlinePlus. Biliary Drainage. U.S. National Library of Medicine; 2025.
  \item Current Medical Diagnosis and Treatment. McGraw-Hill; 2025.
\end{enumerate}
\clearpage
